\subsection{Full Factorial Results}
\label{sec:marl-full-results}

This section reports results from the complete $5 \times 5 \times 5$ factorial experiment: 125 conditions, 30 replications each, 1,000 training episodes per replication. Total computational budget: 3,750 independent training runs executed on an AMD Radeon RX 7900 XTX (25.8\,GB VRAM) using a GPU-vectorized PyTorch implementation with ROCm 6.3.

\subsubsection{Experimental Summary}

Table~\ref{tab:marl-design} summarises the experimental parameters. The factorial design manipulates alignment ($\alpha \in \{-0.8, -0.4, 0, 0.4, 0.8\}$), stakes ($\sigma \in \{0.2, 0.4, 0.6, 0.8, 1.0\}$), and entropy ($\varepsilon \in \{0, 0.25, 0.5, 0.75, 1.0\}$). Each of the 125 conditions was replicated 30 times with independent random seeds using Independent Q-Learning (IQL) agents. Each replication trained for 1,000 episodes of 100 timesteps, with the final 100 episodes used for evaluation metrics.

Across all 3,750 replications, mean reward ranged from $-4.511$ (worst condition: $\alpha = 0.0$, $\sigma = 1.0$, $\varepsilon = 0.75$) to $-0.425$ (best condition: $\alpha = 0.4$, $\sigma = 0.2$, $\varepsilon = 0.00$), a ratio of approximately $10.6\times$ in reward gap magnitude. The grand mean was $-1.724$ with standard deviation $1.252$.

\subsubsection{Hypothesis Tests}

\textbf{H1: Alignment--Friction Inverse Relationship.} Table~\ref{tab:marl-main-effects} reports the factorial ANOVA results. Alignment exhibits a significant main effect ($F = 186.76$, $\eta^2 = 0.0958$, $p < 0.001$). However, the effect is \emph{not} monotonically inverse as H1 predicts. Instead, the alignment--friction relationship is U-shaped: neutral alignment ($\alpha = 0$) produces the worst outcomes (mean reward $-2.486$), while both cooperative ($\alpha = 0.4$: $-1.456$) and adversarial ($\alpha = -0.4$: $-1.453$) alignment yield comparable performance. The linear regression coefficient $\hat{\beta}_\alpha = -0.017$ is near zero, reflecting this symmetry.

This U-shape has a mechanistic interpretation under IQL. When agents have correlated targets ($\alpha > 0$), their independent learning trajectories naturally converge toward compatible policies. When targets are anti-correlated ($\alpha < 0$), agents develop complementary strategies---each agent learns to avoid what others pursue. At neutral alignment ($\alpha = 0$), targets are uncorrelated, providing no structural signal for IQL agents to exploit. The non-stationarity of independent learning is maximally damaging when the reward landscape offers no exploitable alignment structure.

H1 is therefore \emph{partially} supported: alignment matters, and the best outcomes occur at cooperative alignment ($\alpha = 0.4$). But the relationship is non-monotonic, with adversarial alignment outperforming neutrality---a finding that refines the theoretical prediction. The friction framework predicts monotonic increase as $\alpha \to -1$ (via the $1/(1+\alpha)$ divergence), but the simulation reveals that alignment \emph{structure}---whether positive or negative---is more valuable than alignment \emph{direction} for IQL coordination.

\textbf{H2: Stakes--Friction Positive Relationship.} Stakes exhibit the strongest main effect by a substantial margin ($F = 772.55$, $\eta^2 = 0.3965$, $p < 0.001$), confirming H2. The regression coefficient $\hat{\beta}_\sigma = 2.786$ indicates that stakes is the dominant predictor of coordination failure, accounting for approximately 40\% of total variance. Mean reward decreases monotonically from $-0.596$ at $\sigma = 0.2$ to $-2.821$ at $\sigma = 1.0$---a nearly $5\times$ increase in coordination loss. The stakes--friction relationship is approximately linear, consistent with the theoretical prediction of homogeneous-degree-1 scaling in $\sigma$.

\textbf{H3: Entropy--Friction Positive Relationship.} Entropy shows a significant but modest main effect ($F = 15.30$, $\eta^2 = 0.0079$, $p < 0.001$). The regression coefficient $\hat{\beta}_\varepsilon = 0.286$ confirms that observation noise increases friction, supporting H3. Mean reward decreases from $-1.557$ at $\varepsilon = 0$ to $-1.854$ at $\varepsilon = 1.0$, a 19\% increase in coordination loss. The effect is roughly an order of magnitude weaker than stakes, suggesting that in this resource allocation environment, what agents want matters far more than how accurately they perceive the state. The $\alpha \times \varepsilon$ interaction is non-significant ($F = 1.10$, $p = 0.350$), indicating that entropy amplifies friction approximately additively with alignment rather than multiplicatively as the theoretical $(1+\varepsilon)/(1+\alpha)$ form would predict.

\textbf{H4: Friction Function Fit.} Table~\ref{tab:marl-model-comparison} reports model comparison results. The theoretical friction specification M1 ($F = \sigma(1+\varepsilon)/(1+\alpha)$) achieves $R^2 = 0.108$ on condition-level reward gaps, well below the pre-registered threshold of $R^2 > 0.7$. The poor fit is driven by two factors: (1) the U-shaped alpha effect violates the monotonic $1/(1+\alpha)$ assumption, and (2) the dominance of stakes over the other parameters means that collapsing all three into a single composite index loses substantial information.

The independent-effects model M4 ($\Delta R \sim \alpha + \sigma + \varepsilon$) achieves $R^2 = 0.753$ with the lowest AIC ($165.7$), outperforming M1 by $\Delta\text{AIC} = 156.6$---decisive evidence by information-theoretic criteria. Note that M4's near-zero alpha coefficient ($\hat{\beta}_\alpha = -0.017$) is not evidence that alignment is irrelevant; rather, the linear term fails to capture the U-shaped relationship. A quadratic specification $\hat{\beta}_{\alpha^2} \cdot \alpha^2$ would substantially improve even M4.

Critically, all four models agree on the directional predictions for stakes and entropy (H2, H3). The disagreement concerns the alignment functional form and the multiplicative interaction structure. The friction function captures the right variables but not the right shape for alignment.

\subsubsection{Interaction Effects}

Table~\ref{tab:marl-interactions} reports the full interaction structure. The $\alpha \times \sigma$ interaction is significant ($F = 10.24$, $\eta^2 = 0.021$, $p < 0.001$), confirming superadditivity: the worst-performing conditions cluster at neutral alignment \emph{and} high stakes ($\alpha = 0$, $\sigma \geq 0.8$). Figure~\ref{fig:marl-heatmap-alpha-sigma} visualises this interaction as a heatmap of reward gap across the $\alpha \times \sigma$ plane. The U-shaped alpha effect is amplified at high stakes: the performance difference between $\alpha = 0$ and $\alpha = \pm 0.4$ grows with $\sigma$.

The $\sigma \times \varepsilon$ interaction is marginally significant ($F = 1.85$, $\eta^2 = 0.004$, $p = 0.021$), suggesting that entropy has a slightly larger effect at high stakes. The three-way interaction $\alpha \times \sigma \times \varepsilon$ is non-significant ($F = 0.94$, $p = 0.623$), providing limited support for the fully multiplicative structure of the friction function.

\subsubsection{Convergence Analysis: Dynamic Equilibria}

A striking divergence emerges between policy-level and outcome-level convergence. Using a policy-stability criterion (change $< \delta$ over consecutive sampling windows), only \textbf{0.85\%} of replications achieve policy convergence within the training budget. By contrast, using a reward-stability criterion (10\% tolerance of final reward level), \textbf{99.3\%} of replications achieve reward convergence.

This gap is the simulation's strongest confirmation of the friction framework's core theoretical claim: stable coordination outcomes arise through ongoing mutual adaptation, not through convergence to fixed-point agreement. Agents continuously adjust their policies in response to others' behavior---the joint policy never stabilises---yet the aggregate reward outcome reaches a stable attractor. The system occupies a dynamic equilibrium: individual strategies cycle while collective outcomes persist.

Low-stakes conditions ($\sigma \leq 0.4$) achieve reward convergence within 100--200 episodes, while high-stakes conditions ($\sigma \geq 0.8$) require 300--500 episodes. This pattern holds across alignment levels, confirming that convergence speed scales with coordination difficulty rather than alignment structure. Late-training reward stability (standard deviation over the last 250 episodes) is remarkably uniform across conditions: approximately 0.04--0.06 for all alignment levels, with neutral alignment showing marginally higher instability (0.056 vs 0.043--0.050).

\subsubsection{Agent Inequality}

Beyond aggregate performance, we examine how friction affects the \emph{distribution} of rewards across agents. Agent reward variance exhibits a striking pattern across alignment levels:

\begin{center}
\begin{tabular}{@{}lcc@{}}
\toprule
\textbf{Alignment} & \textbf{Agent Reward Variance} & \textbf{Significant Exploitation} \\
\midrule
$\alpha = -0.8$ & 0.110 & 0\% of conditions \\
$\alpha = -0.4$ & 0.652 & 12\% \\
$\alpha = 0.0$ & \textbf{2.841} & 12\% \\
$\alpha = +0.4$ & 0.636 & 8\% \\
$\alpha = +0.8$ & 0.101 & 0\% \\
\bottomrule
\end{tabular}
\end{center}

The neutral alignment condition ($\alpha = 0$) produces agent reward variance \textbf{28 times higher} than the alignment extremes ($\alpha = \pm 0.8$). This is the most striking equity finding: friction---whether cooperative or adversarial---acts as a structural equaliser. At $|\alpha| = 0.8$, zero conditions exhibit statistically significant exploitation (persistent best-vs-worst agent inequality by $t$-test, $p < 0.05$), while moderate friction and neutral conditions show exploitation in 8--12\% of conditions.

The mechanism differs by sign. Under cooperative friction ($\alpha > 0$), correlated preferences create naturally overlapping targets, so all agents benefit from similar resource states. Under adversarial friction ($\alpha < 0$), anti-correlated preferences create predictable opposition that agents can learn to navigate, and competitive pressure punishes exploitation through counteradaptation. At neutral alignment ($\alpha = 0$), uncorrelated preferences create arbitrary asymmetries---some agents happen to want states near the resource mean, granting them structural advantages that persist without friction to counteract.

The Gini coefficient of agent rewards (mean: 0.397, range: 0.252--0.747) confirms this pattern, with inequality peaking at neutral alignment and high stakes. Stakes amplifies the inequality effect monotonically: agent variance increases from 0.108 at $\sigma = 0.2$ to 1.884 at $\sigma = 1.0$, combining multiplicatively with the alignment effect. The worst inequality occurs at $(\alpha = 0, \sigma = 1.0)$: high stakes with no alignment structure.

\subsubsection{Effect Sizes}

The dynamic range of the experimental manipulation is substantial. Comparing the best condition ($\alpha = 0.4$, $\sigma = 0.2$, $\varepsilon = 0$: mean reward $-0.425 \pm 0.144$) against the worst condition ($\alpha = 0.0$, $\sigma = 1.0$, $\varepsilon = 0.75$: mean reward $-4.511 \pm 1.646$), Cohen's $d = 3.49$ for mean reward, representing a very large effect. The ratio of reward gaps ($10.6\times$) exceeds the theoretical prediction for this parameter range, driven by the unexpectedly severe performance at neutral alignment.

For the theoretically predicted extremes---low friction ($\alpha = 0.8$, $\sigma = 0.2$, $\varepsilon = 0$) vs high friction ($\alpha = -0.8$, $\sigma = 1.0$, $\varepsilon = 1.0$)---Cohen's $d = 2.28$, still a large effect. The low-friction condition achieves mean reward $-0.493 \pm 0.245$ compared to $-3.068 \pm 1.580$ for high friction.

\subsubsection{Cross-Implementation Validation}

To verify that results are not an artifact of the GPU-vectorized implementation, we ran a parallel CPU-based factorial experiment using Python multiprocessing (22 workers) with identical hyperparameters and training budget (1,000 episodes). Table~\ref{tab:marl-cross-validation} reports the comparison across 65 overlapping conditions. Despite using independent random seed streams and different floating-point accumulation patterns (batched GPU tensor operations vs sequential CPU processing), the Spearman rank correlation is $\rho = 0.937$ and the Pearson correlation is $r = 0.927$. The GPU implementation produces systematically lower rewards (64/65 conditions, mean $\Delta = -0.459$), attributable to numerical differences between GPU and CPU execution rather than any algorithmic divergence. The strong rank-order agreement confirms that both implementations identify the same condition-level performance structure: the U-shaped alpha effect and stakes dominance are reproduced across independent codepaths.

\subsubsection{Limitations}

Several limitations qualify these findings. First, IQL is a deliberately naive algorithm choice---coordination failure is what we measure, not what we optimise away. More sophisticated algorithms (MADDPG, QMIX, MAPPO) would likely reduce absolute friction levels while preserving the relative ordering across conditions. Second, the parameter grid is coarse: five levels per factor may miss nonlinearities between grid points, particularly near the divergence at $\alpha \to -1$ and in the transition region around $\alpha = 0$.

Third, the U-shaped alignment effect is potentially specific to IQL and the resource allocation environment. With centralised training (MADDPG) or value decomposition (QMIX), adversarial alignment might produce worse outcomes than neutral alignment, restoring the monotonic prediction. The finding that alignment \emph{structure} matters more than alignment \emph{direction} under IQL is theoretically interesting but may not generalise.

Fourth, convergence time is censored at 1,200 episodes, compressing the upper tail and reducing statistical power for this metric. Nearly all conditions hit the censoring bound, rendering convergence time uninformative for distinguishing between moderate- and high-friction conditions. Future work should extend training duration or adopt uncensored convergence criteria.

Finally, the poor performance of the theoretical friction model M1 ($R^2 = 0.108$) suggests that the specific multiplicative functional form $\sigma(1+\varepsilon)/(1+\alpha)$ requires revision. The monotonic $1/(1+\alpha)$ term fails to capture the U-shaped alignment effect, and the multiplicative interaction structure is not supported by the non-significant three-way interaction. A revised specification incorporating $\alpha^2$ or $|\alpha|$ terms may better fit the empirical pattern. However, the qualitative predictions---that coordination failure is structured by alignment, stakes, and entropy---are robustly confirmed.

\subsubsection{Summary}

The MARL factorial experiment confirms that friction is a structured, predictable phenomenon in multi-agent coordination. H2 (stakes--friction positive) and H3 (entropy--friction positive) are strongly supported: stakes dominates the reward gap ($\eta^2 = 0.397$), with entropy contributing a smaller but significant effect ($\eta^2 = 0.008$). H1 (alignment--friction inverse) is partially supported: alignment has a significant effect ($\eta^2 = 0.096$), but the relationship is U-shaped rather than monotonic, with neutral alignment producing the worst outcomes. H4 (friction function fit) is not met at the pre-registered threshold for the single-predictor model ($R^2 = 0.108$), though the independent-effects model achieves $R^2 = 0.753$.

Three key findings emerge. First, friction---operationalised as coordination failure in multi-agent reinforcement learning---is governed primarily by stakes magnitude, with alignment contributing a non-trivial but non-monotonic effect. The theoretical friction function $F = \sigma(1+\varepsilon)/(1+\alpha)$ captures the right variables but not the right shape: the U-shaped alignment effect suggests that alignment \emph{structure} (whether cooperative or adversarial) reduces friction under independent learning, while alignment \emph{absence} (neutrality) maximises it.

Second, friction acts as a structural equaliser: agent reward variance drops 28-fold from neutral alignment to the alignment extremes, with zero statistically significant exploitation at $|\alpha| = 0.8$. This is not a tradeoff between efficiency and equity---structured friction improves \emph{both} aggregate performance and distributional fairness simultaneously.

Third, 99.3\% of conditions achieve reward convergence despite only 0.85\% achieving policy convergence---the strongest empirical confirmation that stable coordination arises through dynamic equilibria rather than fixed-point consensus. This finding directly validates the friction framework's prediction that institutional stability requires structured disagreement, not agreement.

These results motivate a refinement of the friction function to incorporate alignment magnitude $|\alpha|$ or $\alpha^2$ terms alongside alignment direction, capturing the empirical finding that alignment structure matters more than alignment sign.
